\documentclass[11pt]{article}
\usepackage[a4paper,margin=24mm]{geometry}
\usepackage{fontspec}
\usepackage{xcolor}
\usepackage{amsmath,amssymb}
\usepackage{booktabs,longtable,array}
\usepackage[hidelinks]{hyperref}
\setmainfont[Path=fonts/,Extension=.otf,UprightFont=*-Regular,BoldFont=*-Bold]{NotoSerifCJKkr}
\definecolor{reviewercolor}{RGB}{38,82,126}
\definecolor{responsecolor}{RGB}{36,113,78}
\definecolor{revisioncolor}{RGB}{145,82,24}
\setlength{\parindent}{0pt}
\setlength{\parskip}{0.55em}
\setlength{\fboxsep}{7pt}
\setlength{\emergencystretch}{3em}
\renewcommand{\familydefault}{\rmdefault}
\newcommand{\fieldheading}[2]{\par\medskip\noindent{\color{#1}\rule{3pt}{1.25em}}\hspace{0.55em}\textbf{#2}\par\smallskip}
\newenvironment{reviewitem}[1]{\subsection*{#1}}{\par\medskip\hrule\bigskip}
\newenvironment{reviewercomment}{\fieldheading{reviewercolor}{1. 리뷰어 코멘트}\begin{quote}}{\end{quote}}
\newenvironment{authorresponse}{\fieldheading{responsecolor}{2. 저자 답변}}{\par}
\newenvironment{manuscriptrevision}[1]{\fieldheading{revisioncolor}{3. 논문 반영 결과}\textbf{수정 위치:} #1\par\smallskip}{\par}
\newcommand{\manuscripttitle}{연속 인과관계(CoC) 주석의 시간·조건 일관성 검사 프레임워크: 적용 가능성과 근거 추적}
\newcommand{\manuscriptid}{확인 필요 (제공 기록에서 확인되지 않음)}
\newcommand{\revisionround}{확인 필요 (개정 소스 버전: revision-v01)}
\newcommand{\responsedate}{2026년 10월 1일}
% Generated from the final manuscript auxiliary labels.
\expandafter\def\csname mloc@sec:intro\endcsname{1절(1쪽)}
\expandafter\def\csname mloc@sec:related\endcsname{2절(1쪽)}
\expandafter\def\csname mloc@fig:overview-pipeline\endcsname{그림 1(2쪽)}
\expandafter\def\csname mloc@sec:data\endcsname{3절(2쪽)}
\expandafter\def\csname mloc@sec:provenance\endcsname{3.1절(2쪽)}
\expandafter\def\csname mloc@fig:annotation-flow\endcsname{그림 2(2쪽)}
\expandafter\def\csname mloc@sec:selection\endcsname{3.2절(2쪽)}
\expandafter\def\csname mloc@tab:related-position\endcsname{표 1(3쪽)}
\expandafter\def\csname mloc@tab:provenance\endcsname{표 2(3쪽)}
\expandafter\def\csname mloc@sec:artifacts\endcsname{3.3절(3쪽)}
\expandafter\def\csname mloc@sec:problem\endcsname{4절(3쪽)}
\expandafter\def\csname mloc@sec:error-types\endcsname{4.1절(3쪽)}
\expandafter\def\csname mloc@eq:hold-conflict\endcsname{식 1(3쪽)}
\expandafter\def\csname mloc@eq:missing-evidence\endcsname{식 5(3쪽)}
\expandafter\def\csname mloc@sec:semantics\endcsname{4.2절(4쪽)}
\expandafter\def\csname mloc@eq:active-update\endcsname{식 7(4쪽)}
\expandafter\def\csname mloc@sec:issue-definition\endcsname{4.3절(4쪽)}
\expandafter\def\csname mloc@eq:response-contract\endcsname{식 9(4쪽)}
\expandafter\def\csname mloc@sec:method\endcsname{5절(4쪽)}
\expandafter\def\csname mloc@sec:workflow\endcsname{5.1절(4쪽)}
\expandafter\def\csname mloc@sec:extraction\endcsname{5.2절(4쪽)}
\expandafter\def\csname mloc@fig:method-pipeline\endcsname{그림 3(5쪽)}
\expandafter\def\csname mloc@tab:worked-example\endcsname{표 3(5쪽)}
\expandafter\def\csname mloc@sec:models\endcsname{5.3절(5쪽)}
\expandafter\def\csname mloc@lst:condition-step\endcsname{목록 1(5쪽)}
\expandafter\def\csname mloc@tab:transition\endcsname{표 4(5쪽)}
\expandafter\def\csname mloc@sec:queries\endcsname{5.4절(5쪽)}
\expandafter\def\csname mloc@fig:condition-observer\endcsname{그림 4(6쪽)}
\expandafter\def\csname mloc@lst:queries\endcsname{목록 2(6쪽)}
\expandafter\def\csname mloc@sec:trace\endcsname{5.5절(6쪽)}
\expandafter\def\csname mloc@tab:scenario-variants\endcsname{표 5(6쪽)}
\expandafter\def\csname mloc@sec:model-scope\endcsname{5.6절(6쪽)}
\expandafter\def\csname mloc@sec:experiments\endcsname{6절(6쪽)}
\expandafter\def\csname mloc@sec:legacy-test\endcsname{6.1절(6쪽)}
\expandafter\def\csname mloc@sec:scenario-design\endcsname{6.2절(6쪽)}
\expandafter\def\csname mloc@sec:retest-design\endcsname{6.3절(7쪽)}
\expandafter\def\csname mloc@sec:aux-design\endcsname{6.4절(7쪽)}
\expandafter\def\csname mloc@tab:regression\endcsname{표 6(7쪽)}
\expandafter\def\csname mloc@sec:results\endcsname{7절(7쪽)}
\expandafter\def\csname mloc@sec:regression-results\endcsname{7.1절(7쪽)}
\expandafter\def\csname mloc@sec:first-results\endcsname{7.2절(7쪽)}
\expandafter\def\csname mloc@tab:first-results\endcsname{표 7(8쪽)}
\expandafter\def\csname mloc@tab:failures\endcsname{표 8(8쪽)}
\expandafter\def\csname mloc@tab:retest\endcsname{표 9(8쪽)}
\expandafter\def\csname mloc@sec:retest-results\endcsname{7.3절(8쪽)}
\expandafter\def\csname mloc@tab:implementation\endcsname{표 10(8쪽)}
\expandafter\def\csname mloc@tab:agreement\endcsname{표 11(8쪽)}
\expandafter\def\csname mloc@tab:reviewer-distribution\endcsname{표 12(8쪽)}
\expandafter\def\csname mloc@sec:agreement-results\endcsname{7.4절(8쪽)}
\expandafter\def\csname mloc@sec:human-results\endcsname{7.5절(8쪽)}
\expandafter\def\csname mloc@sec:trajectory-results\endcsname{7.6절(9쪽)}
\expandafter\def\csname mloc@sec:discussion\endcsname{8절(9쪽)}
\expandafter\def\csname mloc@sec:threats\endcsname{9절(9쪽)}
\expandafter\def\csname mloc@tab:disagreement-examples\endcsname{표 13(10쪽)}
\expandafter\def\csname mloc@tab:trajectory-cases\endcsname{표 14(10쪽)}
\expandafter\def\csname mloc@sec:conclusion\endcsname{10절(10쪽)}

\newcommand{\mloc}[1]{\csname mloc@#1\endcsname}
\newcommand{\responsefront}[2]{%
\begin{center}{\LARGE\bfseries 심사 의견 답변서\par}\vspace{1em}{\large\bfseries\manuscripttitle\par}\end{center}
\vspace{0.8em}
\begin{tabular}{@{}ll@{}}논문 접수번호: & \manuscriptid \\ 수정 차수: & \revisionround \\ 답변일: & \responsedate\end{tabular}
\section*{답변에 앞서}
편집위원장님과 심사위원님께,

연속 CoC의 검사 범위, 의미 변환과 평가 근거를 구체화할 수 있도록 의견을 주셔서 감사합니다.
이번 개정은 자연 오류 성능 대신 시간·조건 관계의 검사 적용 가능성과 근거 추적으로 주장을 제한했습니다.
이미 수행한 재분석과 통제 시나리오의 최초 결과, 보완 후 동일 사례 재시험을 구분하여 반영했습니다.
사람 재평가나 개발 미사용 독립 시험을 수행했다고 주장하지 않습니다.
개정 원고의 파란색은 원 저장소의 최초 제출본 대비 이번에 추가·수정한 내용을 나타냅니다.
검정색 clean 원고는 동일 내용입니다. 아래 절·쪽·표 번호는 이번 개정 원고를 기준으로 합니다.
접수번호와 실제 심사 수정 차수는 기록이 없어 확인 필요로 남겼습니다.

\section*{심사위원 #1}
\textbf{총평(요지)}\par
\begin{quote}#2\end{quote}}
